%!TEX root = ../template.tex
\typeout{NT FILE abstrac-pt.tex}%

As ontologias fornecem conhecimento partilhado referente a um domínio. Fazem-no através da especificação dos conceitos relevantes, e suas relações, num dado domínio.

Em geral, contêm apenas dados públicos e não sensíveis mas, por vezes, podem ser usadas em conjunto de sistemas de apoio à decisão clínica e repositórios de dados médicos. Nestes sistemas existe manipulação frequente de dados pessoais, de carácter sensível.

Atualmente, com o gradual aumento da dimensão das ontologia, a indústria está evoluir no sentido da utilização de ferramentas colaborativas e \emph{outsourcing} de computações pesadas, recorrendo a infrastruturas \emph{cloud} e outras plataformas de computação partilhada. Estas soluções, embora inevitáveis, apresentam novas questões de segurança.

Nesta tese, investigamos técnicas existentes de computação seguras e propomos a definição de um novo protocolo de Encriptação Simétrica Pesquisável, que suporta \emph{queries} SPARQL expressivas sobre ontologias encriptadas.

\begin{keywords}
Ontologias, Encriptação Pesquisável, Computação Segura, SPARQL
\end{keywords} 



